% Dokumenteinstellungen und Anpassungen

\documentclass[chapterprefix=false, 12pt, a4paper, oneside, parskip=half, listof=totoc, bibliography=totoc, numbers=noendperiod, headinclude, twoside]{scrreprt}

%Zeilenabstand einstellen
\usepackage{setspace}

%leertext einfügen
\usepackage{lipsum}

% Für die Erstellung des Hintergrundbilds / shipout befehl
\usepackage{eso-pic}

%Anpassung der Seitenränder (nach Theisen)
\usepackage[top=40mm,bottom=20mm,left=40mm,right=20mm,headsep = \baselineskip, footskip = \dimexpr2\baselineskip+3mm\relax,]{geometry}

% Representation of Directory Structures
\usepackage{dirtree}

%Unterstützung von Umlauten und anderen Sonderzeichen (UTF-8)
\usepackage{lmodern}
\usepackage[utf8]{inputenc}
\usepackage[T1]{fontenc}
\setlength{\emergencystretch}{1em}

% Schriftart Garamond (nach Theisen)
%\usepackage{ebgaramond}

% Schriftart Palatino (nach Theisen)
\usepackage{palatino}

\usepackage{csquotes}
%\usepackage[square,sort,comma,numbers]{natbib}
\usepackage[
backend=biber,
style=alphabetic,
sorting=nyt
]{biblatex}


%\bibliographystyle{agsm}

\usepackage{tabularx} % in the preamble

%Erlaubt unter anderem Umbrüche captions
\usepackage{caption}

%Stichwortverzeichnis
\usepackage{imakeidx}

%Kompakte Listen
\usepackage{paralist}

%Zitate besser formatieren und darstellen
\usepackage{epigraph}

%Ermöglicht Verknüpfungen innerhalb des Dokumentes (e.g. for PDF), Links werden durch "hidelink" nicht explizit hervorgehoben
\usepackage[hidelinks,german]{hyperref}

\usepackage{pdfpages}

%Glossar, Stichwortverzeichnis (Akronyme werden als eigene Liste aufgeführt)
\usepackage[toc, acronym, shortcuts]{glossaries}

%Anpassung von Kopf- und Fußzeile
%beeinflusst die erste Seite des Kapitels
\usepackage[automark, headsepline, autooneside=false]{scrlayer-scrpage}

%Tweaks für scrbook
\usepackage{scrhack}
%\usepackage{fancyhdr}

%\input{resources/styles/header_footer}
\setlength{\parindent}{0em}

%\renewcommand\chapter{\thispagestyle{plain}}

%Auskommentieren für die Verkleinerung des vertikalen Abstandes eines neuen Kapitels
%\renewcommand*{\chapterheadstartvskip}{\vspace*{.25\baselineskip}}
%\renewcommand*{\chapterheadendvskip}{\vspace*{.25\baselineskip}}


%Verbesserte Darstellung der Buchstaben zueinander
\usepackage[stretch=10]{microtype}

%Deutsche Bezeichnungen für angezeigte Namen (z.B. Inhaltsverzeichnis etc.)
\usepackage[ngerman]{babel}

%Einfachere Zitate
\usepackage{epigraph}

%Unterstützung der H Positionierung (keine automatische Verschiebung eingefügter Elemente)
\usepackage{float}

%Erlaubt Umbrüche innerhalb von Tabellen
\usepackage{tabularx}

%Erlaubt Seitenumbrüche innerhalb von Tabellen
\usepackage{longtable}

%Erlaubt die Darstellung von Sourcecode mit Highlighting
\usepackage{listings}

%Vektorgrafiken
%\usepackage{tikz}

%Grafiken (wie jpg, png, etc.)
\usepackage{graphicx}

%Grafiken von Text umlaufen lassen
\usepackage{wrapfig}

\usepackage[dvipsnames]{xcolor}

% Checkmarks
\usepackage{pifont}

%Pseudocode einfügen
\usepackage{algorithm}
\usepackage{algpseudocode}

%Für lange Tabellen
\usepackage{longtable}
\usepackage{longfigure}
\usepackage{capt-of}
%\newcommand*{\footcite}[1]{\footnote{\citep{#1}}}

% Hintergrundbild einfügen
\newcommand\BackgroundPic{%
\put(0,0){%
\parbox[b][\paperheight]{\paperwidth}{%
\vfill
\centering
\includegraphics[width=\paperwidth,height=\paperheight,%
keepaspectratio]{images/deckblatt.png}%
\vfill
}}}

% Neuer Befehl um Subsubsubkapitel schreiben zu können
\newcommand{\subsubsubsection}[1]{\paragraph{#1}\mbox{}\\}

% Schriftgröße der Fußnote ändern
\renewcommand{\footnotesize}{\normalsize}

%Hilfsbefehl für Glossar UND Abkürzung
\newcommand{\newdefinedabbreviation}[4]{
    \newglossaryentry{#1}{
        text={#2},
        long={#3},
        name={\glsentrylong{#1} (\glsentrytext{#1})},
        first={\glsentryname{#1}},
        firstplural={\glsentrylong{#1}\glspluralsuffix (\glsentryname{#1}\glspluralsuffix )},
        description={#4}
    }
}

\newcommand{\cmark}{\ding{51}}%
\newcommand{\xmark}{\ding{55}}%

\newcommand*{\newdualentry}[5][]{%
  \newglossaryentry{main-#2}{name={#4},%
  text={#3\glsadd{#2}},%
  description={{#5}},%
  #1
  }%
  \newglossaryentry{#2}{
  type=\acronymtype,
  first={#4 (#3)},
  name={#3\glsadd{main-#2}},
  description={\glslink{main-#2}{#4}}
  }%
}

% Erzeugen der Deckblätter
\newcommand{\makeTitles}{

% Entfernen der Seitenzahlen und BackroundPic als Hintergrund nutzen
\pagenumbering{gobble}
\AddToShipoutPicture{\BackgroundPic}

% Variablen für das Deckblatt 1
\gradeType{Bachelor of Science (B.Sc.)}
\germanTitle{Deutscher Titel}
\englishTitle{Englischer Titel}
\authorFirstname{Vorname Autor1}
\authorLastname{Nachname Autor1}
\authorBirthplace{Geburtsort Autor1}
\discipline{Fachbereich}
\courseOfStudies{Studiengang}
\matrikelnr{Matrikelnummer}
\submitDate{{\today}}
\firstExaminer{Prof. Dr. Erster Prüfer}
\secondExaminer{Prof. Dr. Zweiter Prüfer}

% Titelblatt 1 erzeugen
\maketitle
\let\cleardoublepage\clearpage
% Eidesstattliche Erklärung 1 erzeugen
\makeeidesstatt
\let\cleardoublepage\clearpage

% Variablen für das Deckblatt 2
\gradeType{Bachelor of Science (B.Sc.)}
\germanTitle{Deutscher Titel}
\englishTitle{Englischer Titel}
\authorFirstname{Vorname Autor2}
\authorLastname{Nachname Autor2}
\authorBirthplace{Geburtsort Autor2}
\discipline{Fachbereich}
\courseOfStudies{Studiengang}
\matrikelnr{Matrikelnummer}
\submitDate{{\today}}
\firstExaminer{Prof. Dr. Erster Prüfer}
\secondExaminer{Prof. Dr. Zweiter Prüfer}

% Titelblatt 2 erzeugen
\maketitle

% Eidesstattliche Erklärung 2 erzeugen
\makeeidesstatt

\ClearShipoutPicture
}


\graphicspath{ {./images/} }

\setlength{\headheight}{18pt}

%Zusätzliche Farben
\definecolor{darkgreen}{RGB}{0,100,0}
\setlength{\emergencystretch}{15pt}

%Umbenennungen
\renewcommand{\lstlistlistingname}{Quelltextverzeichnis}

\renewcommand{\lstlistingname}{Codeauszug}
\lstset{
	language=Java,
	numbers=left,
	columns=fullflexible,
	aboveskip=5pt,
	belowskip=10pt,
	basicstyle=\small\ttfamily,
	backgroundcolor=\color{black!5},
	commentstyle=\color{darkgreen},
	keywordstyle=\color{blue},
	stringstyle=\color{gray},
	showspaces=false,
	showstringspaces=false,
	showtabs=false,
	xleftmargin=16pt,
	xrightmargin=0pt,
	framesep=5pt,
	framerule=3pt,
	frame=leftline,
	%rulecolor=\color{green},
	tabsize=2,
	breaklines=true,
	breakatwhitespace=true,
	prebreak={\mbox{$\hookleftarrow$}}
}

\definecolor{delim}{RGB}{20,105,176}
\definecolor{numb}{RGB}{106, 109, 32}
\definecolor{string}{rgb}{0.64,0.08,0.08}

% Define JSON language style
\lstdefinelanguage{json}{
    showspaces=false,
    showtabs=false,
    breaklines=true,
    postbreak=\raisebox{0ex}[0ex][0ex]{\ensuremath{\color{gray}\hookrightarrow\space}},
    breakatwhitespace=true,
    basicstyle=\ttfamily\small,
    upquote=true,
    morestring=[b]",
    stringstyle=\color{string},
    literate=
     *{0}{{{\color{numb}0}}}{1}
      {1}{{{\color{numb}1}}}{1}
      {2}{{{\color{numb}2}}}{1}
      {3}{{{\color{numb}3}}}{1}
      {4}{{{\color{numb}4}}}{1}
      {5}{{{\color{numb}5}}}{1}
      {6}{{{\color{numb}6}}}{1}
      {7}{{{\color{numb}7}}}{1}
      {8}{{{\color{numb}8}}}{1}
      {9}{{{\color{numb}9}}}{1}
      {\{}{{{\color{delim}{\{}}}}{1}
      {\}}{{{\color{delim}{\}}}}}{1}
      {[}{{{\color{delim}{[}}}}{1}
      {]}{{{\color{delim}{]}}}}{1},
}

\lstdefinelanguage{Kotlin}{
  comment=[l]{//},
  commentstyle={\color{gray}\ttfamily},
  emph={filter, first, firstOrNull, forEach, lazy, map, mapNotNull, println},
  emphstyle={\color{OrangeRed}},
  identifierstyle=\color{black},
  keywords={!in, !is, abstract, actual, annotation, as, as?, break, by, catch, class, companion, const, constructor, continue, crossinline, data, delegate, do, dynamic, else, enum, expect, external, false, field, file, final, finally, for, fun, get, if, import, in, infix, init, inline, inner, interface, internal, is, lateinit, noinline, null, object, open, operator, out, override, package, param, private, property, protected, public, receiveris, reified, return, return@, sealed, set, setparam, super, suspend, tailrec, this, throw, true, try, typealias, typeof, val, var, vararg, when, where, while},
  keywordstyle={\color{NavyBlue}\bfseries},
  morecomment=[s]{/*}{*/},
  morestring=[b]",
  morestring=[s]{"""*}{*"""},
  ndkeywords={@Deprecated, @JvmField, @JvmName, @JvmOverloads, @JvmStatic, @JvmSynthetic, Array, Byte, Double, Float, Int, Integer, Iterable, Long, Runnable, Short, Any, Unit, Nothing},
  ndkeywordstyle={\color{BurntOrange}\bfseries},
  sensitive=true,
  stringstyle={\color{ForestGreen}\ttfamily},
}

\lstset{literate=%
  {Ö}{{\"O}}1
  {Ä}{{\"A}}1
  {Ü}{{\"U}}1
  {ß}{{\ss}}1
  {ü}{{\"u}}1
  {ä}{{\"a}}1
  {ö}{{\"o}}1
}

%Anpassungen für das Abkürzungsverzeichnis
\newglossarystyle{dottedlocations}{%
	\renewcommand*{\glossaryentryfield}[5]{%
		\item[\glsentryitem{##1}\glstarget{##1}{##2}] \emph{##3}%
		\unskip\leaders\hbox to 2.9mm{\hss.}\hfill##5}%
	\renewcommand*{\glsgroupskip}{}%
}


\newglossaryentry{claude}
{
        name=Claude,
        description={Ist ein \acs{llm} von Anthropic PBC}
}

\newglossaryentry{chatgpt}
{
        name=ChatGPT,
        description={Ist ein \acs{llm} von OpenAi}
}

\newglossaryentry{gemini}
{
        name=Gemini,
        description={Ist ein \acs{llm} von Google}
}

\newglossaryentry{llm}
{
        name=large language model,
        description={Ein Large Language Model ist eine Art KI-Modell welche mit großen Datenmengen trainiert wird um aus ihnen zu lernen und somit neue Inhalte generieren zu können. \textbf{QUELLE}}
}

\newglossaryentry{openSource}
{
        name=Open Source,
        description={Open Source bedeutet Quelloffen. Software die Open Source ist, wird von der Öffentlichen Community entwickelt und mit einer frei verfügbaren Lizenz für jeden Zugänglich gemacht.}
}

\newglossaryentry{many2many}
{
        name=Many-To-Many,
        description={Beschreibt eine relationale Datenbankbeziehung. Viele-zu-Viele Beziehung}
}


%\newdualentry{api}{API}{Application Programming Interface}{An Application Programming Interface (API)}

%\newdualentry{crud}{CRUD}{Create, Read, Update, Delete}{CRUD ist ein Akronym für die Vier Grundlegenden Operationen von persistenten Speichern (Create, Read, Update, Delete)}

%\newdualentry{seo}{SEO}{Search Engine Optimization}{Search Engine Optimization wird betrieben, um in Suchmaschinen möglichst weit oben in der Liste zu erscheinen}
\newacronym{seo}{SEO}{Search Engine Optimization}
\newacronym{crud}{CRUD}{Create, Read, Update, Delete}
\newacronym{api}{API}{Application Programming Interface}
\newacronym{http}{HTTP}{Hypertext Transfer Protocol}
\newacronym{ki}{KI}{Künstliche Intelligenz}
\newacronym{rest}{REST}{Representational State Transfer}
\newacronym{json}{JSON}{JavaScript Object Notation}
\newacronym{odm}{ODM}{Object-document Mapping}
\newacronym{orm}{ORM}{Object-relational Mapping}
\newacronym{cors}{CORS}{Cross-Origin Resource Sharing}
\newacronym{di}{DI}{Dependency Injection}
\newacronym{ux}{UX}{User Experience}
\newacronym{ui}{UI}{User Interface}
\newacronym{ide}{IDE}{Integrated Development Environment}
\newacronym{dbms}{DBMS}{Database Management System}
\newacronym{acid}{ACID}{Atomicity, Consistency, Isolation
and Durability}
\newacronym{el}{EL}{Expression Language}
\newacronym{dao}{DAO}{Data Access Object}
\newacronym{sql}{SQL}{Structured Query Language}
\newacronym{url}{URL}{Uniform Resource Locator}
\newacronym{er}{ER}{Entity-Relationship}
\newacronym{mvvm}{MVVM}{Model View ViewModel}
\newacronym{etl}{ETL}{Extract, Transform, Load}
\newacronym{mvc}{MVC}{Model-View-Controller}
\newacronym{gif}{GIF}{Graphics Interchange Format}
\newacronym{html}{HTML}{Hypertext Markup Language}
\newacronym{dsl}{DSL}{Domain Specific Language}
\newacronym{ai}{AI}{Artificial Intelligence}
\newacronym{ifis}{if(is)}{Institut für Internet-Sicherheit}
%\newacronym{llm}{LLM}{Large Language Model}

\include{listings}

% Kapitel des Inhaltsverzeichnis zum Inhaltsverzeichnis hinzufügen
%\setuptoc{toc}{totoc}

\RedeclareSectionCommand[%
  beforeskip=0pt,
  afterskip=1.5\baselineskip
]{chapter}

\RedeclareSectionCommand[%
  beforeskip=-1.0\baselineskip plus -0.75\baselineskip minus -0.5\baselineskip,
  afterskip=0.5\baselineskip plus 0.25\baselineskip minus 0.125\baselineskip
]{section}

\RedeclareSectionCommand[%
  beforeskip=-0.5\baselineskip plus -0.375\baselineskip minus -0.25\baselineskip,
  afterskip=0.25\baselineskip plus 0.125\baselineskip minus 0.0625\baselineskip
]{subsection}

\RedeclareSectionCommand[%
  beforeskip=0pt,
  afterskip=2pt
]{subsubsection}

\makeglossaries
\addbibresource{quellen.bib}

\raggedbottom
\begin{document}

% Header stylen
\clearpairofpagestyles
\ihead{\raisebox{0pt}[\ht\strutbox]{\includegraphics[width=2cm,height=1cm]{wh_logo.png}}}
\chead{\leftmark}
\ohead{\pagemark}

%footer stylen
\ifoot*{}
\cfoot*{}
\ofoot*{}

\makeatletter
\newcommand*{\germanTitle}[1]{\gdef\@germanTitle{#1}}
\newcommand*{\englishTitle}[1]{\gdef\@englishTitle{#1}}
\newcommand*{\gradeType}[1]{\gdef\@gradeType{#1}}
\newcommand*{\firstExaminer}[1]{\gdef\@firstExaminer{#1}}
\newcommand*{\secondExaminer}[1]{\gdef\@secondExaminer{#1}}
\newcommand*{\matrikelnr}[1]{\gdef\@matrikelnr{#1}}
\newcommand*{\authorBirthplace}[1]{\gdef\@authorBirthplace{#1}}
\newcommand*{\submitDate}[1]{\gdef\@submitDate{#1}}
\newcommand*{\discipline}[1]{\gdef\@discipline{#1}}
\newcommand*{\courseOfStudies}[1]{\gdef\@courseOfStudies{#1}}
\newcommand*{\authorLastname}[1]{\gdef\@authorLastname{#1}}
\newcommand*{\authorFirstname}[1]{\gdef\@authorFirstname{#1}}

\renewcommand*{\maketitle}{
	\begin{titlepage}
		\newgeometry{left=3.8cm,right=2.5cm,top=4.6cm,bottom=2.5cm}
			\begingroup
				\fontsize{44pt}{46pt}\selectfont
				{\bfseries Bachelorarbeit}
			\endgroup

			\vskip 1.44cm

			\begingroup
			\fontsize{8pt}{6pt}\selectfont
			Titel der Arbeit // Title of Thesis
			\endgroup

			\vskip -0.02cm

			\begingroup
			\fontsize{16pt}{16pt}\selectfont
			{\bfseries \@germanTitle\par}
			{\bfseries \@englishTitle\par}
			\endgroup
			\vskip -0.1cm

			\noindent\rule{15cm}{0.4pt}
			
			\vskip 0.05cm

			\begingroup
			\fontsize{8pt}{6pt}\selectfont
			Akademischer Abschlussgrad: Grad, Fachrichtung (Abkürzung) // Degree
			\endgroup

			\vskip -0.15cm

			\begingroup
			\fontsize{12pt}{14pt}\selectfont
				{\@gradeType}
			\endgroup
			\vskip -0.15cm

			\noindent\rule{15cm}{0.4pt}
			
			\vskip 0.05cm

			\begingroup
			\fontsize{8pt}{6pt}\selectfont
			Autorenname, Geburtsort // Name, Place of Birth
			\endgroup

			\vskip -0.15cm

			\begingroup
			\fontsize{12pt}{14pt}\selectfont
				{\@authorFirstname} {\@authorLastname}, {\@authorBirthplace}
			\endgroup
			\vskip -0.15cm

			\noindent\rule{15cm}{0.4pt}
			
			\vskip 0.05cm

			\begingroup
			\fontsize{8pt}{6pt}\selectfont
			Studiengang // Course of Study
			\endgroup

			\vskip -0.15cm

			\begingroup
			\fontsize{12pt}{14pt}\selectfont
				{\@courseOfStudies}
			\endgroup
			\vskip -0.15cm

			\noindent\rule{15cm}{0.4pt}
			
			\vskip 0.05cm

			\begingroup
			\fontsize{8pt}{6pt}\selectfont
			Fachbereich // Department
			\endgroup

			\vskip -0.15cm

			\begingroup
			\fontsize{12pt}{14pt}\selectfont
				{\@discipline}
			\endgroup
			\vskip -0.15cm

			\noindent\rule{15cm}{0.4pt}
			
			\vskip 0.05cm

			\begingroup
			\fontsize{8pt}{6pt}\selectfont
			Erstprüferin/Erstprüfer // First Examiner
			\endgroup

			\vskip -0.15cm

			\begingroup
			\fontsize{12pt}{14pt}\selectfont
				{\@firstExaminer}
			\endgroup
			\vskip -0.15cm

			\noindent\rule{15cm}{0.4pt}
			
			\vskip 0.05cm

			\begingroup
			\fontsize{8pt}{6pt}\selectfont
			Zweitprüferin/Zweitprüfer // Second Examiner
			\endgroup

			\vskip -0.15cm

			\begingroup
			\fontsize{12pt}{14pt}\selectfont
				{\@secondExaminer}
			\endgroup
			\vskip -0.15cm

			\noindent\rule{15cm}{0.4pt}
			
			\vskip 0.05cm

			\begingroup
			\fontsize{8pt}{6pt}\selectfont
			Abgabedatum // Date of Submission
			\endgroup

			\vskip -0.15cm

			\begingroup
			\fontsize{12pt}{14pt}\selectfont
				{8. Mai 2024}
			\endgroup
			\vskip -0.15cm

			\noindent\rule{15cm}{0.4pt}
		\restoregeometry
	\end{titlepage}
}
\makeatother

\makeatletter

\newcommand*{\place}[1]{\gdef\@place{#1}}

\newcommand*{\makeeidesstatt}{
	\begin{titlepage}
		\newgeometry{left=3.8cm,right=2.5cm,top=8cm,bottom=2.5cm}
			\begingroup
			\fontsize{18pt}{20pt}\selectfont
			{\bfseries Eidesstattliche Versicherung}
			\endgroup

			\vskip 0.8cm

			\begingroup
			\fontsize{12pt}{14pt}\selectfont
				{\@authorLastname}, {\@authorFirstname}
			\endgroup

			\vskip -0.5cm

			\noindent\rule{15cm}{0.4pt}

			\vskip -0.4cm

			\begingroup
			\fontsize{8pt}{6pt}\selectfont
			Name, Vorname // Name, First Name
			\endgroup

			\vskip 0.6cm
			\begingroup
			\fontsize{10.5pt}{11.5pt}\selectfont
			Ich versichere hiermit an Eides statt, dass ich die vorliegende Abschlussarbeit mit dem Titel
			\endgroup

			\begingroup
			\fontsize{16pt}{18pt}\selectfont
			{\bfseries \@germanTitle}
			\endgroup

			\begingroup
			\fontsize{10.5pt}{11.5pt}\selectfont
			selbstständig und ohne unzulässige fremde Hilfe erbracht habe. Ich habe keine anderen als die angegebenen Quellen und Hilfsmittel benutzt sowie wörtliche und sinngemä{\ss}e Zitate kenntlich gemacht. Die Arbeit hat in gleicher oder ähnlicher Form noch keiner Prüfungsbehörde vorgelegen.
			\endgroup

			\vskip 0.8cm

			{
			\fontsize{12pt}{14pt}\selectfont
			{\@place}, den
			}

			\vskip -0.5cm

			\noindent\rule{15cm}{0.4pt}

			\vskip -0.3cm

			\begingroup
			\fontsize{8pt}{6pt}\selectfont
			Ort, Datum, Unterschrift // Place, Date, Signature
			\endgroup


		\restoregeometry
	\end{titlepage}
}
\makeatother


\setcounter{secnumdepth}{4}
\setcounter{tocdepth}{4}
\newcounter{savepagearab}
\newcounter{savepageroman}

\title{Bachelorarbeit}
\place{Gelsenkirchen}

% Deckblätter einfügen
\makeTitles

\begin{onehalfspace}
 
% Abstract und Vorworte
\thispagestyle{empty}

\chapter*{Abstract}
Meistens auf Englisch

\chapter*{Zusammenfassung}

Deutsche Kurzfassung/Zusammenfassung
\newpage

% Inhaltsverzeichnis einfügen

% Seitennummerierung in Römisch für die Verzeichnisse
\pagenumbering{Roman}
% Header bei Kapiteln einfügen
\renewcommand*{\chapterpagestyle}{scrheadings}

\setcounter{page}{7}
\tableofcontents
\newpage

%\renewcommand{\listfigurename}{Darstellungsverzeichnis}
%\renewcommand{\figurename}{Darst.}
% Abbildungsverzeichnis einfügen
\listoffigures

\listoftables

\lstlistoflistings

% Abkürzungsverzeichnis einfügen
%\thispagestyle{empty}
\printglossary[type=\acronymtype,title=Abkürzungsverzeichnis, toctitle=Abkürzungsverzeichnis]
\newpage

% Aktuelle Seite speichern
\setcounter{savepageroman}{\arabic{page}}
\pagenumbering{arabic}


\chapter{Einleitung}
 
\section{Zielsetzung / Problemstellung}

\section{Forschungsstand}

\section{Ausgangssituation}


% Theorie Einführung
\chapter{Kapitelname}

\section{Unterkapitel}
\chapter{Kapitelname}

\section{Unterkapitel}
\chapter{Kapitel mit Bild}

\begin{figure} [ht]
        \begin{center}
        \includegraphics[width=\textwidth]{wh_logo}
        \caption{Beispielbild}
        \label{fig:ER-Diagramm}
        \end{center}
\end{figure}

\section{Unterkapitel}
\chapter{Kapitelname}

\section{Unterkapitel}
\chapter{Fazit}

\section{Zusammenfassung der Ergebnisse}


\section{Retrospektive}


\section{Ausblick}


\newpage 
%\pagenumbering{Roman}
%\setcounter{page}{\thesavepageroman}
\printglossary 

% Bibliografie im .bib Format einfügen
\printbibliography[title=Literaturverzeichnis]

\appendix
\chapter{Anhang}

\section{Anhang 1}

\subsection{Dokument oder Text}

\newpage

\end{onehalfspace}



%\bibliography{quellen}\addcontentsline{toc}{section}{\protect{}Literaturverzeichnis}%
\end{document}

